\documentclass[border=10pt]{standalone}
\usepackage{fontspec}
\setmainfont{Noto Serif}[Ligatures=TeXOff]
\XeTeXgenerateactualtext=1
\newfontfamily\bengalifont{Noto Serif Bengali}[Script=Bengali, NFSSFamily=bnfam]
\usepackage[Latin,Bengali]{ucharclasses}
% group-free font switching (as in header-common.tex): robust inside braces, cells and quotes
\makeatletter
\def\BN@saved{\rmdefault}
\def\BN@bn{bnfam}
\newcommand\BNin{\edef\BN@tmp{\f@family}\ifx\BN@tmp\BN@bn\else\let\BN@saved\BN@tmp\fi\fontfamily{bnfam}\selectfont}
\newcommand\BNout{\fontfamily{\BN@saved}\selectfont}
\makeatother
\setTransitionsFor{Bengali}{\BNin}{\BNout}
% Typographic quotes and dashes are not in a ucharclasses class (class 0), so a Bangla word right
% after “ or ‘ would not trigger the Bangla font. Treat them as Latin so that every switch fires.
\XeTeXcharclass"2018=\XeTeXcharclass`A
\XeTeXcharclass"2019=\XeTeXcharclass`A
\XeTeXcharclass"201C=\XeTeXcharclass`A
\XeTeXcharclass"201D=\XeTeXcharclass`A
\XeTeXcharclass"2013=\XeTeXcharclass`A
\XeTeXcharclass"2014=\XeTeXcharclass`A
\XeTeXcharclass"2026=\XeTeXcharclass`A
\usepackage{newunicodechar}
\newunicodechar{।}{{\bengalifont\char"0964}}
\newunicodechar{॥}{{\bengalifont\char"0965}}

\usepackage[table]{xcolor}
\usepackage{booktabs,tabularx,array}
\definecolor{bapilrule}{HTML}{1F2937}
\definecolor{bapilstripe}{HTML}{F3F4F6}
\renewcommand{\arraystretch}{1.28}
\setlength{\tabcolsep}{5pt}

\begin{document}
\begin{minipage}{17cm}
{\large\bfseries সারণি ৩০।}\ {\large বাংলা নীতি ফাউন্ডেশন মডেল: প্রমাণ-নিয়ন্ত্রিত ধাপক্রম}\par\vspace{7pt}
\small
\rowcolors{2}{bapilstripe}{white}
\begin{tabularx}{\linewidth}{@{}>{\hsize=0.619\hsize\raggedright\arraybackslash}X>{\hsize=1.143\hsize\raggedright\arraybackslash}X>{\hsize=1.190\hsize\raggedright\arraybackslash}X>{\hsize=0.810\hsize\raggedright\arraybackslash}X>{\hsize=1.238\hsize\raggedright\arraybackslash}X@{}}
\arrayrulecolor{bapilrule}\toprule
\rowcolor{white}\textbf{ধাপ} & \textbf{প্রবেশের মানদণ্ড} & \textbf{প্রধান কার্যক্রম} & \textbf{উৎপাদ} & \textbf{প্রস্থান বা গ্রহণের মানদণ্ড} \\
\midrule
\textbf{পর্যায় ০ক} শব্দভিত্তিক ও মেটাডেটা অনুসন্ধান & মেটাডেটাসহ নিয়ন্ত্রিত P0 কর্পাস & দ্বিভাষিক সম্প্রসারণসহ BM25; নথির ধরন, প্রকাশক সংস্থা ও বছরভিত্তিক ফিল্টার; বিরতি & BanglaPolicyQA v0-তে ভিত্তিমান ফলাফল & ভিত্তিমান প্রকাশিত; পরবর্তী প্রতিটি ধাপের স্বচ্ছ বিকল্প হিসেবে বহাল \\
\textbf{পর্যায় ০খ} ঘন ভেক্টরভিত্তিক উদ্ধার ও পুনঃক্রমবিন্যাস & পর্যায় ০ক-এর ভিত্তিমান; উন্মুক্ত বহুভাষিক এমবেডিং ও পুনঃক্রমবিন্যাস মডেলের সংক্ষিপ্ত তালিকা (সারণি ২৯) & মিশ্র উদ্ধার; পুনঃক্রমবিন্যাস; ভাষা ও নথির ধরনভিত্তিক তুলনামূলক বিশ্লেষণ & মিশ্র কনফিগারেশন; মূল্যায়ন প্রতিবেদন & Recall@10 বা nDCG@10 আস্থা-ব্যবধানের চেয়ে বেশি বাড়লে এবং কোনো নিরাপত্তা-মাপকাঠি খারাপ না হলেই গৃহীত \\
\textbf{পর্যায় ০গ} বিদ্যমান মডেলে উত্তর & উদ্ধার তার সীমা অর্জন করে; উদ্ধৃতি-বিচারের প্রটোকল প্রস্তুত (সংযুক্তি দ) & নিষ্কাশনমূলক উত্তর; বাক্যভিত্তিক উদ্ধৃতি-যাচাইসহ উন্মুক্ত মডেলে ঐচ্ছিক পুনর্লিখন & উত্তর সেবা; উদ্ধৃতির ফলাফল & উদ্ধৃতি-নির্ভুলতা ও অসমর্থিত দাবির সীমা অর্জিত (সারণি ২৮); অন্যথায় শুধু নিষ্কাশনমূলক উত্তর \\
\textbf{পর্যায় ১} লক্ষ্যভিত্তিক অভিযোজন & পর্যায় ০-এর পর অবশিষ্ট ত্রুটি শ্রেণিভিত্তিকভাবে নথিভুক্ত এবং উদ্ধার বা ডেটা নয়, মডেলের কারণে বলে চিহ্নিত; বাছাইকৃত ডেটা; কম্পিউটিং; ERRC পর্যালোচনা & প্যারামিটার-দক্ষ অভিযোজন (LoRA); বাছাইকৃত জোড়ায় নির্দেশনা-টিউনিং; রেড-টিমিং & অভিযোজিত মডেল; মডেল কার্ড; ব্যয়-প্রতিবেদন & সম্মত মাপকাঠিতে পর্যায় ০গ-এর তুলনায় আস্থা-ব্যবধানের চেয়ে বেশি উন্নতি; সব নিরাপত্তা-সীমা অর্জিত \\
\textbf{পর্যায় ২} নিবেদিত ডোমেইন মডেল & অগ্রাধিকারমূলক কাজে পর্যায় ১ অপর্যাপ্ত; পর্যাপ্ত নিয়ন্ত্রিত বাংলা টেক্সট; জাতীয় কম্পিউটিং; ব্যয়-উপযোগিতার যুক্তি অনুমোদিত; জাতীয় বাংলা LLM উদ্যোগের সঙ্গে সমন্বয় & নিয়ন্ত্রিত বাংলা টেক্সটে বিদ্যমান উন্মুক্ত মডেলের অব্যাহত প্রাক-প্রশিক্ষণ; স্বাধীন নিরীক্ষা & ডোমেইন মডেল প্রকাশ; নিরীক্ষা প্রতিবেদন & স্বাধীন নিরীক্ষা উত্তীর্ণ; স্টিয়ারিং বোর্ড ও ERRC-এর প্রকাশ-সিদ্ধান্ত। শূন্য থেকে নতুন মডেল প্রশিক্ষণ আওতার বাইরে (অনুচ্ছেদ ৯.৭) \\
\bottomrule
\end{tabularx}
\par\vspace{3pt}{\scriptsize\color{gray} উৎস: BAPIL টেকনিক্যাল ব্লুপ্রিন্ট, সংস্করণ ১.০ (২০২৬)। CC BY 4.0।}
\end{minipage}
\end{document}
